\section{環境法・循環・持続可能性}

\subsection{平成23年度 問5-1}
\label{q:h23-reg-basic-5-1}

\noindent\textbf{原典：} sources/H23/basic.pdf、PDF物理15頁。

\textbf{公式正答：} 1（single）。

\par\textit{原典本文の校正と独自解説は未作成。分類のみの登録項目。}

\subsection{平成23年度 問5-2}
\label{q:h23-reg-basic-5-2}

\noindent\textbf{原典：} sources/H23/basic.pdf、PDF物理16頁。

\textbf{公式正答：} 2（single）。

\par\textit{原典本文の校正と独自解説は未作成。分類のみの登録項目。}

\subsection{平成24年度 問5-1}
\label{q:h24-reg-basic-5-1}

\noindent\textbf{原典：} sources/H24/basic.pdf、PDF物理15頁。

\textbf{公式正答：} 3（single）。

\par\textit{原典本文の校正と独自解説は未作成。分類のみの登録項目。}

\subsection{平成25年度 問5-3}
\label{q:h25-reg-basic-5-3}

\noindent\textbf{原典：} sources/H25/basic.pdf、PDF物理20頁。

\textbf{公式正答：} 4（single）。

\par\textit{原典本文の校正と独自解説は未作成。分類のみの登録項目。}

\subsection{平成25年度 問5-4}
\label{q:h25-reg-basic-5-4}

\noindent\textbf{原典：} sources/H25/basic.pdf、PDF物理21頁。

\textbf{公式正答：} 2（single）。

\par\textit{原典本文の校正と独自解説は未作成。分類のみの登録項目。}

\subsection{平成26年度 問5-2}
\label{q:h26-reg-basic-5-2}

\noindent\textbf{原典：} sources/H26/basic.pdf、PDF物理16頁。

\textbf{公式正答：} 5（single）。

\par\textit{原典本文の校正と独自解説は未作成。分類のみの登録項目。}

\subsection{平成27年度 問5-1}
\label{q:h27-reg-basic-5-1}

\noindent\textbf{原典：} sources/H27/basic.pdf、PDF物理15頁。

\textbf{公式正答：} 2（single）。

\par\textit{原典本文の校正と独自解説は未作成。分類のみの登録項目。}

\subsection{平成27年度 問5-2}
\label{q:h27-reg-basic-5-2}

\noindent\textbf{原典：} sources/H27/basic.pdf、PDF物理15頁。

\textbf{公式正答：} 1（single）。

\par\textit{原典本文の校正と独自解説は未作成。分類のみの登録項目。}

\subsection{平成28年度 問5-1}
\label{q:h28-reg-basic-5-1}

\noindent\textbf{原典：} sources/H28/basic.pdf、PDF物理21頁。

\textbf{公式正答：} 5（single）。

\par\textit{原典本文の校正と独自解説は未作成。分類のみの登録項目。}

\subsection{平成29年度 問5-1}
\label{q:h29-reg-basic-5-1}

\noindent\textbf{原典：} sources/H29/basic.pdf、PDF物理18頁。

\textbf{公式正答：} 1（single）。

\par\textit{原典本文の校正と独自解説は未作成。分類のみの登録項目。}

\subsection{平成30年度 問5-1}
\label{q:h30-reg-basic-5-1}

\noindent\textbf{原典：} sources/H30/basic.pdf、PDF物理18頁。

\textbf{公式正答：} 2（single）。

\par\textit{原典本文の校正と独自解説は未作成。分類のみの登録項目。}

\subsection{平成30年度 問5-2}
\label{q:h30-reg-basic-5-2}

\noindent\textbf{原典：} sources/H30/basic.pdf、PDF物理18頁。

\textbf{公式正答：} 3（single）。

\par\textit{原典本文の校正と独自解説は未作成。分類のみの登録項目。}

\subsection{令和元年度 問5-1}
\label{q:r01-reg-basic-5-1}

\noindent\textbf{原典：} sources/R01/basic.pdf、PDF物理19頁。

\textbf{公式正答：} 3（single）。

\par\textit{原典本文の校正と独自解説は未作成。分類のみの登録項目。}

\subsection{令和元年度 問5-2}
\label{q:r01-reg-basic-5-2}

\noindent\textbf{原典：} sources/R01/basic.pdf、PDF物理20頁。

\textbf{公式正答：} 2（single）。

\par\textit{原典本文の校正と独自解説は未作成。分類のみの登録項目。}

\subsection{令和元年度再試験 問5-1}
\label{q:r01-ret-basic-5-1}

\noindent\textbf{原典：} sources/R01\_retest/basic.pdf、PDF物理17頁。

\textbf{公式正答：} 1（single）。

\par\textit{原典本文の校正と独自解説は未作成。分類のみの登録項目。}

\subsection{令和元年度再試験 問5-2}
\label{q:r01-ret-basic-5-2}

\noindent\textbf{原典：} sources/R01\_retest/basic.pdf、PDF物理18頁。

\textbf{公式正答：} 1（single）。

\par\textit{原典本文の校正と独自解説は未作成。分類のみの登録項目。}

\subsection{令和元年度再試験 問5-4}
\label{q:r01-ret-basic-5-4}

\noindent\textbf{原典：} sources/R01\_retest/basic.pdf、PDF物理19頁。

\textbf{公式正答：} 3（single）。

\par\textit{原典本文の校正と独自解説は未作成。分類のみの登録項目。}

\subsection{令和2年度 問5-1}
\label{q:r02-reg-basic-5-1}

\noindent\textbf{原典：} sources/R02/basic.pdf、PDF物理19頁。

\textbf{公式正答：} 2（single）。

\par\textit{原典本文の校正と独自解説は未作成。分類のみの登録項目。}

\subsection{令和3年度 問5-2}
\label{q:r03-reg-basic-5-2}

\noindent\textbf{原典：} sources/R03/basic.pdf、PDF物理22頁。

\textbf{公式正答：} 5（single）。

\par\textit{原典本文の校正と独自解説は未作成。分類のみの登録項目。}

\subsection{令和3年度 問5-3}
\label{q:r03-reg-basic-5-3}

\noindent\textbf{原典：} sources/R03/basic.pdf、PDF物理23頁。

\textbf{公式正答：} 1（single）。

\par\textit{原典本文の校正と独自解説は未作成。分類のみの登録項目。}

\subsection{令和4年度 問5-2}
\label{q:r04-reg-basic-5-2}

\noindent\textbf{原典：} sources/R04/basic.pdf、PDF物理17頁。

\textbf{公式正答：} 1（single）。

\par\textit{原典本文の校正と独自解説は未作成。分類のみの登録項目。}

\subsection{令和6年度 問5-1}
\label{q:r06-reg-basic-5-1}

\noindent\textbf{原典：} sources/R06/basic.pdf、PDF物理17頁。

\textbf{公式正答：} 1（single）。

\par\textit{原典本文の校正と独自解説は未作成。分類のみの登録項目。}

\subsection{令和6年度 問5-2}
\label{q:r06-reg-basic-5-2}

\noindent\textbf{原典：} sources/R06/basic.pdf、PDF物理18頁。

\textbf{公式正答：} 5（single）。

\par\textit{原典本文の校正と独自解説は未作成。分類のみの登録項目。}

\subsection{令和6年度 問5-3}
\label{q:r06-reg-basic-5-3}

\noindent\textbf{原典：} sources/R06/basic.pdf、PDF物理19頁。

\textbf{公式正答：} 2（single）。

\par\textit{原典本文の校正と独自解説は未作成。分類のみの登録項目。}

\subsection{令和6年度 問5-4}
\label{q:r06-reg-basic-5-4}

\noindent\textbf{原典：} sources/R06/basic.pdf、PDF物理19頁。

\textbf{公式正答：} 4（single）。

\par\textit{原典本文の校正と独自解説は未作成。分類のみの登録項目。}

\subsection{令和7年度 問5-1}
\label{q:r07-reg-basic-5-1}

\noindent\textbf{原典：} sources/R07/basic.pdf、PDF物理15頁。

\textbf{公式正答：} 5（single）。

\par\textit{原典本文の校正と独自解説は未作成。分類のみの登録項目。}

\subsection{令和7年度 問5-2}
\label{q:r07-reg-basic-5-2}

\noindent\textbf{原典：} sources/R07/basic.pdf、PDF物理16頁。

\textbf{公式正答：} 3（single）。

\par\textit{原典本文の校正と独自解説は未作成。分類のみの登録項目。}

\subsection{令和7年度 問5-4}
\label{q:r07-reg-basic-5-4}

\noindent\textbf{原典：} sources/R07/basic.pdf、PDF物理18頁。

\textbf{公式正答：} 2（single）。

\par\textit{原典本文の校正と独自解説は未作成。分類のみの登録項目。}

\section{科学技術史}

\subsection{平成24年度 問5-5}
\label{q:h24-reg-basic-5-5}

\noindent\textbf{原典：} sources/H24/basic.pdf、PDF物理17頁。

\textbf{公式正答：} 4（single）。

\par\textit{原典本文の校正と独自解説は未作成。分類のみの登録項目。}

\subsection{平成29年度 問5-5}
\label{q:h29-reg-basic-5-5}

\noindent\textbf{原典：} sources/H29/basic.pdf、PDF物理19頁。

\textbf{公式正答：} 2（single）。

\par\textit{原典本文の校正と独自解説は未作成。分類のみの登録項目。}

\subsection{令和6年度 問5-5}
\label{q:r06-reg-basic-5-5}

\noindent\textbf{原典：} sources/R06/basic.pdf、PDF物理20頁。

\textbf{公式正答：} 3（single）。

\par\textit{原典本文の校正と独自解説は未作成。分類のみの登録項目。}

\section{知的財産}

\subsection{平成27年度 問5-5}
\label{q:h27-reg-basic-5-5}

\noindent\textbf{原典：} sources/H27/basic.pdf、PDF物理17頁。

\textbf{公式正答：} 3（single）。

\par\textit{原典本文の校正と独自解説は未作成。分類のみの登録項目。}

\subsection{令和元年度 問5-6}
\label{q:r01-reg-basic-5-6}

\noindent\textbf{原典：} sources/R01/basic.pdf、PDF物理22頁。

\textbf{公式正答：} 5（single）。

\par\textit{原典本文の校正と独自解説は未作成。分類のみの登録項目。}

\section{安全・リスク評価}

\subsection{平成25年度 問5-1}
\label{q:h25-reg-basic-5-1}

\noindent\textbf{原典：} sources/H25/basic.pdf、PDF物理19頁。

\textbf{公式正答：} 5（single）。

\par\textit{原典本文の校正と独自解説は未作成。分類のみの登録項目。}

\subsection{平成30年度 問5-3}
\label{q:h30-reg-basic-5-3}

\noindent\textbf{原典：} sources/H30/basic.pdf、PDF物理19頁。

\textbf{公式正答：} 3（single）。

\par\textit{原典本文の校正と独自解説は未作成。分類のみの登録項目。}

\subsection{令和元年度再試験 問5-6}
\label{q:r01-ret-basic-5-6}

\noindent\textbf{原典：} sources/R01\_retest/basic.pdf、PDF物理20頁。

\textbf{公式正答：} 3（single）。

\par\textit{原典本文の校正と独自解説は未作成。分類のみの登録項目。}

\subsection{令和4年度 問5-3}
\label{q:r04-reg-basic-5-3}

\noindent\textbf{原典：} sources/R04/basic.pdf、PDF物理18頁。

\textbf{公式正答：} 1（single）。

\par\textit{原典本文の校正と独自解説は未作成。分類のみの登録項目。}

\subsection{令和4年度 問5-5}
\label{q:r04-reg-basic-5-5}

\noindent\textbf{原典：} sources/R04/basic.pdf、PDF物理20頁。

\textbf{公式正答：} 5（single）。

\par\textit{原典本文の校正と独自解説は未作成。分類のみの登録項目。}

\section{研究・情報倫理}

\subsection{平成25年度 問5-5}
\label{q:h25-reg-basic-5-5}

\noindent\textbf{原典：} sources/H25/basic.pdf、PDF物理22頁。

\textbf{公式正答：} 2（single）。

\par\textit{原典本文の校正と独自解説は未作成。分類のみの登録項目。}

\subsection{平成28年度 問5-5}
\label{q:h28-reg-basic-5-5}

\noindent\textbf{原典：} sources/H28/basic.pdf、PDF物理24頁。

\textbf{公式正答：} 5（single）。

\par\textit{原典本文の校正と独自解説は未作成。分類のみの登録項目。}

\section{科学技術史・政策の経緯}

\subsection{平成23年度 問5-5}
\label{q:h23-reg-basic-5-5}

\noindent\textbf{原典：} sources/H23/basic.pdf、PDF物理17頁。

\textbf{公式正答：} 4（single）。

\par\textit{原典本文の校正と独自解説は未作成。分類のみの登録項目。}

\subsection{平成25年度 問5-6}
\label{q:h25-reg-basic-5-6}

\noindent\textbf{原典：} sources/H25/basic.pdf、PDF物理22頁。

\textbf{公式正答：} 1（single）。

\par\textit{原典本文の校正と独自解説は未作成。分類のみの登録項目。}

\subsection{平成26年度 問5-6}
\label{q:h26-reg-basic-5-6}

\noindent\textbf{原典：} sources/H26/basic.pdf、PDF物理18頁。

\textbf{公式正答：} 4（single）。

\par\textit{原典本文の校正と独自解説は未作成。分類のみの登録項目。}

\subsection{平成27年度 問5-6}
\label{q:h27-reg-basic-5-6}

\noindent\textbf{原典：} sources/H27/basic.pdf、PDF物理18頁。

\textbf{公式正答：} 4（single）。

\par\textit{原典本文の校正と独自解説は未作成。分類のみの登録項目。}

\subsection{平成28年度 問5-6}
\label{q:h28-reg-basic-5-6}

\noindent\textbf{原典：} sources/H28/basic.pdf、PDF物理24頁。

\textbf{公式正答：} 4（single）。

\par\textit{原典本文の校正と独自解説は未作成。分類のみの登録項目。}

\subsection{平成29年度 問5-6}
\label{q:h29-reg-basic-5-6}

\noindent\textbf{原典：} sources/H29/basic.pdf、PDF物理20頁。

\textbf{公式正答：} 5（single）。

\par\textit{原典本文の校正と独自解説は未作成。分類のみの登録項目。}

\subsection{平成30年度 問5-5}
\label{q:h30-reg-basic-5-5}

\noindent\textbf{原典：} sources/H30/basic.pdf、PDF物理20頁。

\textbf{公式正答：} 3（single）。

\par\textit{原典本文の校正と独自解説は未作成。分類のみの登録項目。}

\subsection{令和元年度再試験 問5-5}
\label{q:r01-ret-basic-5-5}

\noindent\textbf{原典：} sources/R01\_retest/basic.pdf、PDF物理19頁。

\textbf{公式正答：} 4（single）。

\par\textit{原典本文の校正と独自解説は未作成。分類のみの登録項目。}

\subsection{令和2年度 問5-6}
\label{q:r02-reg-basic-5-6}

\noindent\textbf{原典：} sources/R02/basic.pdf、PDF物理22頁。

\textbf{公式正答：} 1（single）。

\par\textit{原典本文の校正と独自解説は未作成。分類のみの登録項目。}

\subsection{令和3年度 問5-5}
\label{q:r03-reg-basic-5-5}

\noindent\textbf{原典：} sources/R03/basic.pdf、PDF物理25頁。

\textbf{公式正答：} 4（single）。

\par\textit{原典本文の校正と独自解説は未作成。分類のみの登録項目。}

\subsection{令和3年度 問5-6}
\label{q:r03-reg-basic-5-6}

\noindent\textbf{原典：} sources/R03/basic.pdf、PDF物理26頁。

\textbf{公式正答：} 4（single）。

\par\textit{原典本文の校正と独自解説は未作成。分類のみの登録項目。}

\subsection{令和4年度 問5-6}
\label{q:r04-reg-basic-5-6}

\noindent\textbf{原典：} sources/R04/basic.pdf、PDF物理20頁。

\textbf{公式正答：} 3（single）。

\par\textit{原典本文の校正と独自解説は未作成。分類のみの登録項目。}

\subsection{令和5年度 問5-5}
\label{q:r05-reg-basic-5-5}

\noindent\textbf{原典：} sources/R05/basic.pdf、PDF物理18頁。

\textbf{公式正答：} 1（single）。

\par\textit{原典本文の校正と独自解説は未作成。分類のみの登録項目。}

\subsection{令和6年度 問5-6}
\label{q:r06-reg-basic-5-6}

\noindent\textbf{原典：} sources/R06/basic.pdf、PDF物理20頁。

\textbf{公式正答：} 2（single）。

\par\textit{原典本文の校正と独自解説は未作成。分類のみの登録項目。}

\subsection{令和7年度 問5-5}
\label{q:r07-reg-basic-5-5}

\noindent\textbf{原典：} sources/R07/basic.pdf、PDF物理18頁。

\textbf{公式正答：} 2（single）。

\par\textit{原典本文の校正と独自解説は未作成。分類のみの登録項目。}

\subsection{令和7年度 問5-6}
\label{q:r07-reg-basic-5-6}

\noindent\textbf{原典：} sources/R07/basic.pdf、PDF物理19頁。

\textbf{公式正答：} 4（single）。

\par\textit{原典本文の校正と独自解説は未作成。分類のみの登録項目。}

\section{独立正規変数の和差と標本平均}

\subsection{平成23年度 問5-4}
\label{q:h23-reg-basic-5-4}

\noindent\textbf{原典：} sources/H23/basic.pdf、PDF物理17頁。

\textbf{公式正答：} 3（single）。

\paragraph{問題の要旨} 独立な製品質量の標本平均と総質量について、平均と標準偏差を求める。

\Needspace{4\baselineskip}\paragraph{確認した条件・記述（独自要約）}\mbox{}\par

\noindent\textbf{判断方向：} 最も適切な記述又は組合せを選ぶ\par

\noindent\textbf{原典照合：} 14原典画像で16問の条件、数値、単位、表の列と全5候補を確認。\par

\noindent\textbf{モデル：} 独立性、正規性、一定発生率などの前提を個別に保持する。\par

\paragraph{問われる核心} 独立な製品質量の標本平均と総質量について、平均と標準偏差を求める。

\paragraph{必要知識} 独立変数の線形結合、平均の線形性、分散への係数の二乗、正規分布の閉性。

\paragraph{考え方・正答を選ぶ理由（査読待ち初稿）} 各質量を独立正規変数とする。Aの標本平均は \(\bar A=(A_1+\cdots+A_4)/4\) で、期待値80g、分散 \(4(4^2)/4^2=4\,\mathrm{g^2}\) である。総質量は \(T=\sum_{i=1}^4A_i+\sum_{i=1}^4B_i\) なので期待値 \(4(80)+4(120)=800\,\mathrm g\)、分散 \(4(4^2)+4(3^2)=100\,\mathrm{g^2}\) となる。標準偏差はそれぞれ2gと10gで、3を選ぶ。平均を取る操作と、四つの値を合計する操作を混同しない。

\Needspace{4\baselineskip}\paragraph{選択肢の検討（査読待ち初稿）}\mbox{}\par

\noindent\textbf{1：不適切} 平均20・偏差4・合計400・偏差10。Aの標本平均の期待値は80gであり、標本数で平均を割らない。総質量の期待値も800gである。\par

\noindent\textbf{2：不適切} 平均80・偏差4・合計800・偏差7。平均の偏差は4gではなく2gである。合計の偏差も標準偏差をそのまま足した7gではなく10gである。\par

\noindent\textbf{3：適切} 平均80・偏差2・合計800・偏差10。Aの4個平均の偏差は4/2=\allowbreak{}2g、合計の分散は4×16+4×9=\allowbreak{}100g\textsuperscript{2}なので10gである。\par

\noindent\textbf{4：不適切} 平均20・偏差2・合計400・偏差7。平均と総質量の期待値がともに違い、合計の偏差も違う。\par

\noindent\textbf{5：不適切} 平均80・偏差2・合計800・偏差7。平均部分は合うが、合計の標準偏差は7gではない。独立成分の分散を足して最後に平方根を取る。\par

\paragraph{再利用する考え方} 変数の意味と単位を先に定め、平均・分散・確率・判定を順に分ける。品質の要求、工程の監視、ロットの受入れを区別し、成立条件のない一般化を避ける。

\paragraph{共通理論} \hyperref[theory:quality:variance]{第\ref*{theory:quality:variance}節}, \hyperref[theory:quality:normal]{第\ref*{theory:quality:normal}節}, \hyperref[theory:quality:sample_mean]{第\ref*{theory:quality:sample_mean}節}

\par 参考：\cite{stat_nist_normal,stat_psu_linear}。

\paragraph{関連問題} \hyperref[q:r06-reg-basic-1-6]{R06-REG-BASIC-1-6}

\par 難易度 3（校準初稿）、優先度未評価。

\section{工程の品質保証と標準化}

\subsection{平成24年度 問5-4}
\label{q:h24-reg-basic-5-4}

\noindent\textbf{原典：} sources/H24/basic.pdf、PDF物理17頁。

\textbf{公式正答：} 2（single）。

\paragraph{問題の要旨} 最終検査と上流の品質保証を分け、プロセスの標準化と改善の関係を読む。

\Needspace{4\baselineskip}\paragraph{確認した条件・記述（独自要約）}\mbox{}\par

\noindent\textbf{判断方向：} 最も適切な記述又は組合せを選ぶ\par

\noindent\textbf{原典照合：} 14原典画像で16問の条件、数値、単位、表の列と全5候補を確認。\par

\noindent\textbf{モデル：} 独立性、正規性、一定発生率などの前提を個別に保持する。\par

\paragraph{問われる核心} 最終検査と上流の品質保証を分け、プロセスの標準化と改善の関係を読む。

\paragraph{必要知識} 検査と予防、上流工程、標準化、PDCA、プロセスと成果。

\paragraph{考え方・正答を選ぶ理由（査読待ち初稿）} 完成品を検査して不適合品を隔離することと、不適合を生む工程を上流から改善することを分ける。設問の文脈ではアが検査、イが源流管理、ウが標準化、エがプロセスであり2となる。標準を設定して実施し、結果を確認して更新する循環が品質保証を支える。標準化は永久に方法を固定することではなく、改善結果を次の標準へ反映することを含む。

\Needspace{4\baselineskip}\paragraph{選択肢の検討（査読待ち初稿）}\mbox{}\par

\noindent\textbf{1：不適切} 検査・三現管理・基準化・プロセス。現場・現物・現実という三現の確認と、上流で不適合原因を抑える源流管理は同じ説明ではない。\par

\noindent\textbf{2：適切} 検査・源流管理・標準化・プロセス。要求適合を判定する検査、上流からの源流管理、方法を定める標準化、プロセスを管理する流れでつながる。\par

\noindent\textbf{3：不適切} 検査・三現管理・基準化・アウトカム。三現管理・基準化の対応に加え、ここで管理する対象は一連のプロセスである。\par

\noindent\textbf{4：不適切} 評価・源流管理・標準化・アウトカム。評価は広い概念だが本問の要求適合判定は検査である。末尾も成果だけを意味するアウトカムではなくプロセスである。\par

\noindent\textbf{5：不適切} 評価・三現管理・標準化・プロセス。要求適合判定と上流の改善に対応する語が検査・源流管理になっていない。\par

\paragraph{再利用する考え方} 変数の意味と単位を先に定め、平均・分散・確率・判定を順に分ける。品質の要求、工程の監視、ロットの受入れを区別し、成立条件のない一般化を避ける。

\paragraph{共通理論} \hyperref[theory:quality:assurance]{第\ref*{theory:quality:assurance}節}, \hyperref[theory:quality:pdca]{第\ref*{theory:quality:pdca}節}

\par 参考：\cite{quality_process_jsqc,quality_iso_process}。

\paragraph{関連問題} \hyperref[q:r03-reg-basic-1-3]{R03-REG-BASIC-1-3}

\par 難易度 2（校準初稿）、優先度未評価。

\section{未確定}

\subsection{平成23年度 問5-3}
\label{q:h23-reg-basic-5-3}

\noindent\textbf{原典：} sources/H23/basic.pdf、PDF物理16頁。

\textbf{公式正答：} 2（single）。

\par\textit{原典本文の校正と独自解説は未作成。分類のみの登録項目。}

\subsection{平成24年度 問5-2}
\label{q:h24-reg-basic-5-2}

\noindent\textbf{原典：} sources/H24/basic.pdf、PDF物理15頁。

\textbf{公式正答：} 4（single）。

\par\textit{原典本文の校正と独自解説は未作成。分類のみの登録項目。}

\subsection{平成24年度 問5-3}
\label{q:h24-reg-basic-5-3}

\noindent\textbf{原典：} sources/H24/basic.pdf、PDF物理16頁。

\textbf{公式正答：} 2（single）。

\par\textit{原典本文の校正と独自解説は未作成。分類のみの登録項目。}

\subsection{平成25年度 問5-2}
\label{q:h25-reg-basic-5-2}

\noindent\textbf{原典：} sources/H25/basic.pdf、PDF物理19頁。

\textbf{公式正答：} 2（single）。

\par\textit{原典本文の校正と独自解説は未作成。分類のみの登録項目。}

\subsection{平成26年度 問5-1}
\label{q:h26-reg-basic-5-1}

\noindent\textbf{原典：} sources/H26/basic.pdf、PDF物理16頁。

\textbf{公式正答：} 1（single）。

\par\textit{原典本文の校正と独自解説は未作成。分類のみの登録項目。}

\subsection{平成26年度 問5-3}
\label{q:h26-reg-basic-5-3}

\noindent\textbf{原典：} sources/H26/basic.pdf、PDF物理17頁。

\textbf{公式正答：} 5（single）。

\par\textit{原典本文の校正と独自解説は未作成。分類のみの登録項目。}

\subsection{平成26年度 問5-4}
\label{q:h26-reg-basic-5-4}

\noindent\textbf{原典：} sources/H26/basic.pdf、PDF物理17頁。

\textbf{公式正答：} 4（single）。

\par\textit{原典本文の校正と独自解説は未作成。分類のみの登録項目。}

\subsection{平成26年度 問5-5}
\label{q:h26-reg-basic-5-5}

\noindent\textbf{原典：} sources/H26/basic.pdf、PDF物理17頁。

\textbf{公式正答：} 3（single）。

\par\textit{原典本文の校正と独自解説は未作成。分類のみの登録項目。}

\subsection{平成27年度 問5-3}
\label{q:h27-reg-basic-5-3}

\noindent\textbf{原典：} sources/H27/basic.pdf、PDF物理16頁。

\textbf{公式正答：} 5（single）。

\par\textit{原典本文の校正と独自解説は未作成。分類のみの登録項目。}

\subsection{平成27年度 問5-4}
\label{q:h27-reg-basic-5-4}

\noindent\textbf{原典：} sources/H27/basic.pdf、PDF物理17頁。

\textbf{公式正答：} 5（single）。

\par\textit{原典本文の校正と独自解説は未作成。分類のみの登録項目。}

\subsection{平成28年度 問5-2}
\label{q:h28-reg-basic-5-2}

\noindent\textbf{原典：} sources/H28/basic.pdf、PDF物理22頁。

\textbf{公式正答：} 5（single）。

\par\textit{原典本文の校正と独自解説は未作成。分類のみの登録項目。}

\subsection{平成28年度 問5-3}
\label{q:h28-reg-basic-5-3}

\noindent\textbf{原典：} sources/H28/basic.pdf、PDF物理22頁。

\textbf{公式正答：} 3（single）。

\par\textit{原典本文の校正と独自解説は未作成。分類のみの登録項目。}

\subsection{平成28年度 問5-4}
\label{q:h28-reg-basic-5-4}

\noindent\textbf{原典：} sources/H28/basic.pdf、PDF物理23頁。

\textbf{公式正答：} 1（single）。

\par\textit{原典本文の校正と独自解説は未作成。分類のみの登録項目。}

\subsection{平成29年度 問5-2}
\label{q:h29-reg-basic-5-2}

\noindent\textbf{原典：} sources/H29/basic.pdf、PDF物理18頁。

\textbf{公式正答：} 1（single）。

\par\textit{原典本文の校正と独自解説は未作成。分類のみの登録項目。}

\subsection{平成29年度 問5-3}
\label{q:h29-reg-basic-5-3}

\noindent\textbf{原典：} sources/H29/basic.pdf、PDF物理19頁。

\textbf{公式正答：} 4（single）。

\par\textit{原典本文の校正と独自解説は未作成。分類のみの登録項目。}

\subsection{平成29年度 問5-4}
\label{q:h29-reg-basic-5-4}

\noindent\textbf{原典：} sources/H29/basic.pdf、PDF物理19頁。

\textbf{公式正答：} 2（single）。

\par\textit{原典本文の校正と独自解説は未作成。分類のみの登録項目。}

\subsection{平成30年度 問5-4}
\label{q:h30-reg-basic-5-4}

\noindent\textbf{原典：} sources/H30/basic.pdf、PDF物理19頁。

\textbf{公式正答：} 4（single）。

\par\textit{原典本文の校正と独自解説は未作成。分類のみの登録項目。}

\subsection{平成30年度 問5-6}
\label{q:h30-reg-basic-5-6}

\noindent\textbf{原典：} sources/H30/basic.pdf、PDF物理20頁。

\textbf{公式正答：} 5（single）。

\par\textit{原典本文の校正と独自解説は未作成。分類のみの登録項目。}

\subsection{令和元年度 問5-3}
\label{q:r01-reg-basic-5-3}

\noindent\textbf{原典：} sources/R01/basic.pdf、PDF物理20頁。

\textbf{公式正答：} 3（single）。

\par\textit{原典本文の校正と独自解説は未作成。分類のみの登録項目。}

\subsection{令和元年度 問5-4}
\label{q:r01-reg-basic-5-4}

\noindent\textbf{原典：} sources/R01/basic.pdf、PDF物理21頁。

\textbf{公式正答：} 4（single）。

\par\textit{原典本文の校正と独自解説は未作成。分類のみの登録項目。}

\subsection{令和元年度 問5-5}
\label{q:r01-reg-basic-5-5}

\noindent\textbf{原典：} sources/R01/basic.pdf、PDF物理21頁。

\textbf{公式正答：} 2（single）。

\par\textit{原典本文の校正と独自解説は未作成。分類のみの登録項目。}

\subsection{令和元年度再試験 問5-3}
\label{q:r01-ret-basic-5-3}

\noindent\textbf{原典：} sources/R01\_retest/basic.pdf、PDF物理18頁。

\textbf{公式正答：} 3（single）。

\par\textit{原典本文の校正と独自解説は未作成。分類のみの登録項目。}

\subsection{令和2年度 問5-2}
\label{q:r02-reg-basic-5-2}

\noindent\textbf{原典：} sources/R02/basic.pdf、PDF物理20頁。

\textbf{公式正答：} 3（single）。

\par\textit{原典本文の校正と独自解説は未作成。分類のみの登録項目。}

\subsection{令和2年度 問5-3}
\label{q:r02-reg-basic-5-3}

\noindent\textbf{原典：} sources/R02/basic.pdf、PDF物理20頁。

\textbf{公式正答：} 4（single）。

\par\textit{原典本文の校正と独自解説は未作成。分類のみの登録項目。}

\subsection{令和2年度 問5-4}
\label{q:r02-reg-basic-5-4}

\noindent\textbf{原典：} sources/R02/basic.pdf、PDF物理21頁。

\textbf{公式正答：} 1（single）。

\par\textit{原典本文の校正と独自解説は未作成。分類のみの登録項目。}

\subsection{令和2年度 問5-5}
\label{q:r02-reg-basic-5-5}

\noindent\textbf{原典：} sources/R02/basic.pdf、PDF物理21頁。

\textbf{公式正答：} 4|3（multiple\_accepted）。

\par\textit{原典本文の校正と独自解説は未作成。分類のみの登録項目。}

\subsection{令和3年度 問5-1}
\label{q:r03-reg-basic-5-1}

\noindent\textbf{原典：} sources/R03/basic.pdf、PDF物理21頁。

\textbf{公式正答：} 2（single）。

\par\textit{原典本文の校正と独自解説は未作成。分類のみの登録項目。}

\subsection{令和3年度 問5-4}
\label{q:r03-reg-basic-5-4}

\noindent\textbf{原典：} sources/R03/basic.pdf、PDF物理24頁。

\textbf{公式正答：} 2（single）。

\par\textit{原典本文の校正と独自解説は未作成。分類のみの登録項目。}

\subsection{令和4年度 問5-1}
\label{q:r04-reg-basic-5-1}

\noindent\textbf{原典：} sources/R04/basic.pdf、PDF物理17頁。

\textbf{公式正答：} 2（single）。

\par\textit{原典本文の校正と独自解説は未作成。分類のみの登録項目。}

\subsection{令和4年度 問5-4}
\label{q:r04-reg-basic-5-4}

\noindent\textbf{原典：} sources/R04/basic.pdf、PDF物理19頁。

\textbf{公式正答：} 5（single）。

\par\textit{原典本文の校正と独自解説は未作成。分類のみの登録項目。}

\subsection{令和5年度 問5-1}
\label{q:r05-reg-basic-5-1}

\noindent\textbf{原典：} sources/R05/basic.pdf、PDF物理16頁。

\textbf{公式正答：} 3（single）。

\par\textit{原典本文の校正と独自解説は未作成。分類のみの登録項目。}

\subsection{令和5年度 問5-2}
\label{q:r05-reg-basic-5-2}

\noindent\textbf{原典：} sources/R05/basic.pdf、PDF物理16頁。

\textbf{公式正答：} 5（single）。

\par\textit{原典本文の校正と独自解説は未作成。分類のみの登録項目。}

\subsection{令和5年度 問5-3}
\label{q:r05-reg-basic-5-3}

\noindent\textbf{原典：} sources/R05/basic.pdf、PDF物理17頁。

\textbf{公式正答：} 1（single）。

\par\textit{原典本文の校正と独自解説は未作成。分類のみの登録項目。}

\subsection{令和5年度 問5-4}
\label{q:r05-reg-basic-5-4}

\noindent\textbf{原典：} sources/R05/basic.pdf、PDF物理17頁。

\textbf{公式正答：} 2（single）。

\par\textit{原典本文の校正と独自解説は未作成。分類のみの登録項目。}

\subsection{令和5年度 問5-6}
\label{q:r05-reg-basic-5-6}

\noindent\textbf{原典：} sources/R05/basic.pdf、PDF物理18頁。

\textbf{公式正答：} 5（single）。

\par\textit{原典本文の校正と独自解説は未作成。分類のみの登録項目。}

\subsection{令和7年度 問5-3}
\label{q:r07-reg-basic-5-3}

\noindent\textbf{原典：} sources/R07/basic.pdf、PDF物理17頁。

\textbf{公式正答：} 4（single）。

\par\textit{原典本文の校正と独自解説は未作成。分類のみの登録項目。}
